\section{Contracts, Realizations, and Repair Reports}
\label{sec:contracts}

\paragraph{Finite interfaces.}
Let \(A\) be a type of contract atoms and \(L\) a type of implementation
levers, both with decidable equality.  The active contract interface is an
explicit finite set \(I \subseteq A\).  An implementation block map
\(\beta : A \to \mathcal{P}_{\mathrm{fin}}(L)\) assigns each atom a finite set
of levers.  The active lever interface is
\[
  \Lambda_I := \bigcup_{a \in I} \beta(a).
\]
For any finite atom set \(T\), write
\[
  \betaSet(T) := \bigcup_{a \in T}\beta(a),
  \qquad\text{so that}\qquad
  \Lambda_I=\betaSet(I).
\]
The active blocks are assumed pairwise disjoint in the theorem statements that
compare block-aligned and grouped repairs.  Nonemptiness is used where a
grouping inverse is required.

\paragraph{Retained-set predicates.}
The reference and implementation residuals are represented abstractly by
\[
  \SatPsi : \mathcal{P}_{\mathrm{fin}}(A) \to \mathsf{Prop},
  \qquad
  \SatPhi : \mathcal{P}_{\mathrm{fin}}(L) \to \mathsf{Prop}.
\]
\(\SatPsi(T)\) means that the contract residual retaining \(T\) is feasible;
\(\SatPhi(K)\) has the analogous implementation meaning.  All definitions use
the retained-set convention.  A deletion is expressed by a complement inside
\(I\) or \(\Lambda_I\), never by a second deletion-based satisfiability
predicate.

\(\SatPsi\) and \(\SatPhi\) are abstract feasibility predicates, not CNF
predicates.  Retained sets are the primitive objects: deleting \(G\) from the
contract interface is represented by \(\SatPsi(I\setminus G)\), while deleting
\(R\) from the implementation interface is represented by
\(\SatPhi(\Lambda_I\setminus R)\).  Keeping complements inside their declared
active domains avoids mixing deletion sets with the retained residuals on
which feasibility is evaluated.

\begin{figure}[t]
  \centering
  \fbox{\parbox{0.91\textwidth}{\centering
    \textbf{Figure placeholder: two-axis framework}\\[0.5em]
    Truth / certificate axis: reference residual
    \(\leftrightarrow\) implementation residual\\[0.4em]
    Repair / report axis: raw lever deletion
    \(\longrightarrow\) grouped contract deletion}}
  \caption{Residual legitimacy and repair reportability are distinct axes.}
  \label{fig:two-axis}
\end{figure}

\begin{definition}[Inclusion-minimality]
For a predicate \(P\) on finite sets,
\[
  \mathsf{InclusionMinimal}(P,X)
  \quad\Longleftrightarrow\quad
  P(X) \ \land\
  \forall Y,\ P(Y) \to Y \subseteq X \to X \subseteq Y.
\]
\end{definition}

\begin{definition}[Raw repair]
A lever deletion \(R\) is raw-feasible when
\[
  R \subseteq \Lambda_I
  \quad\text{and}\quad
  \SatPhi(\Lambda_I \setminus R).
\]
\(\RawRepair(R)\) holds when \(R\) is inclusion-minimal among raw-feasible
deletions.
\end{definition}

\begin{definition}[Touch-any grouping]
The reported atoms touched by a lever deletion are
\[
  \groupTouchAny(R)
  :=
  \{a \in I \mid (R \cap \beta(a)) \neq \varnothing\}.
\]
\end{definition}

\begin{definition}[Contract and grouped repairs]
A contract deletion \(G\) is contract-feasible when
\[
  G \subseteq I
  \quad\text{and}\quad
  \SatPsi(I \setminus G).
\]
\(\ContractRepair(G)\) is its inclusion-minimal form.  A set \(G\) is in the
raw grouped image when some \(\RawRepair(R)\) satisfies
\(\groupTouchAny(R)=G\).  Finally, \(\GroupedRepair(G)\) is an
inclusion-minimal element of that grouped image.
\end{definition}

\begin{definition}[Contract conditions]
The Lean development uses the following conditions.
\begin{align*}
\ResidualFaithfulness
&:\Longleftrightarrow
  \forall T \subseteq I,\
  \SatPhi\!\left(\bigcup_{a\in T}\beta(a)\right)
  \leftrightarrow \SatPsi(T),\\
\RepairAtomicity
&:\Longleftrightarrow
  \forall a\in I,\ |\beta(a)|=1,\\
\GroupSoundness
&:\Longleftrightarrow
  \forall R\subseteq\Lambda_I,\
  \SatPhi(\Lambda_I\setminus R)
  \to \SatPsi(I\setminus\groupTouchAny(R)),\\
\PsiDeletionMonotonicity
&:\Longleftrightarrow
  \forall T'\subseteq T\subseteq I,\
  \SatPsi(T)\to\SatPsi(T').
\end{align*}
\end{definition}

\paragraph{Paper-to-Lean correspondence.}
The core objects above map directly to the declarations in
\code{Defs.lean}:
\begin{description}
  \item[\(\betaSet\)] \code{betaSet};
  \item[\(\Lambda_I\)] \code{LambdaI};
  \item[raw feasible deletion] \code{RawFeasible};
  \item[raw repair] \code{RawRepair};
  \item[touch-any grouping] \code{groupTouchAny};
  \item[contract feasible deletion] \code{ContractFeasible};
  \item[contract repair] \code{ContractRepair};
  \item[grouped repair] \code{GroupedRepair};
  \item[residual faithfulness] \code{ResidualFaithfulness};
  \item[repair atomicity] \code{RepairAtomicity};
  \item[group soundness] \code{GroupSoundness}; and
  \item[\(\Psi\)-deletion monotonicity] \code{PsiDeletionMonotonicity}.
\end{description}
This displayed list is a declaration map rather than an additional
main-text table.  Appendix~\ref{app:lean} gives the corresponding theorem and
lemma map.

These definitions do not assign CNF or solver semantics to \(\SatPsi\) or
\(\SatPhi\).  Such semantics belong to an application-side truth/certificate
contract.
